\documentclass[11pt,a4paper]{article}
\usepackage[italian]{babel}
\usepackage[margin=2.5cm]{geometry}
\usepackage{parskip}
\usepackage{amsmath, amssymb, amsthm}
\usepackage{booktabs}
\usepackage{array}
\usepackage{xcolor}
\usepackage{needspace}
\usepackage{enumitem}
\usepackage{listings}
\usepackage{tikz}
\usepackage[most]{tcolorbox}
\usepackage[hidelinks]{hyperref}
\usetikzlibrary{decorations.pathreplacing, shapes.geometric, arrows.meta, positioning, automata}

% Niente vedove/orfane: i paragrafi non lasciano righe isolate a fine/inizio pagina
\widowpenalty=10000
\clubpenalty=10000

% Elenchi più compatti
\setlist{itemsep=2pt, parsep=0pt, topsep=4pt, partopsep=0pt}

% --- Riquadri con bordo nero, senza numero (si spezzano tra due pagine) ---
% Uso: \begin{definizione} ... \end{definizione}
%      \begin{teorema}[Nome] ... \end{teorema}  (il nome è facoltativo)
\tcbset{nota/.style={enhanced jigsaw, breakable, colback=white,
  colframe=black, boxrule=0.5pt, sharp corners}}
\NewTColorBox{definizione}{}{nota,
  before upper={\textbf{Definizione.}\ }}
\NewTColorBox{teorema}{o}{nota,
  before upper={\textbf{Teorema\IfValueT{#1}{ (#1)}.}\ }}
\NewTColorBox{proposizione}{o}{nota,
  before upper={\textbf{Proposizione\IfValueT{#1}{ (#1)}.}\ }}
\NewTColorBox{proprieta}{o}{nota,
  before upper={\textbf{Proprietà\IfValueT{#1}{ (#1)}.}\ }}
\NewTColorBox{assioma}{o}{nota, boxrule=1pt,
  before upper={\textbf{Assioma\IfValueT{#1}{ (#1)}.}\ }}

% --- Ambienti semplici (senza riquadro) ---
% Il titolo sta in cima al blocco, anche se il contenuto inizia con un riquadro o
% con codice; e non resta mai da solo in fondo a una pagina.
% Uso: \begin{esempio}[Titolo facoltativo] ... \end{esempio}
\newcommand{\newrunin}[2]{%
  \NewDocumentEnvironment{#1}{o}{%
    \par\Needspace{4\baselineskip}\addvspace{\medskipamount}\noindent
    \textbf{#2}\IfValueT{##1}{ (##1)}\textbf{.}\ \ignorespaces}%
  {\par\addvspace{\medskipamount}}}
\newrunin{esempio}{Esempio}
\newrunin{esercizio}{Esercizio}
\NewTColorBox{osservazione}{}{nota,
  before upper={\textbf{Osservazione.}\ }}

% V e F nelle tavole di verità
\newcommand{\V}{\textsf{V}}
\newcommand{\F}{\textsf{F}}

% --- Codice C ---
% Uso: \begin{codice} ... \end{codice}      codice C numerato
%      \begin{comando} ... \end{comando}    comandi da terminale
%      \begin{risultato} ... \end{risultato} output di un programma
%      \lstinline|printf("ciao");|           codice dentro al testo
\lstdefinestyle{c}{
  language=C,
  basicstyle=\ttfamily\small,
  keywordstyle=\color{blue}\bfseries,
  commentstyle=\color{gray}\itshape,
  stringstyle=\color{red!70!black},
  numbers=left,
  numberstyle=\tiny\color{gray},
  numbersep=8pt,
  columns=flexible,
  keepspaces=true,
  breaklines=true,
  showstringspaces=false,
  tabsize=4,
  morekeywords={include, define},
  % lettere accentate nei commenti
  literate={à}{{\`a}}1 {è}{{\`e}}1 {é}{{\'e}}1 {ì}{{\`i}}1
           {ò}{{\`o}}1 {ù}{{\`u}}1 {È}{{\`E}}1 {À}{{\`A}}1
}
\lstset{style=c}
\newtcblisting{codice}{listing only, nota, left=6mm,
  listing options={style=c}}
\newtcblisting{comando}{listing only, nota,
  listing options={basicstyle=\ttfamily\small, numbers=none, language={},
    columns=flexible, keepspaces=true, breaklines=true}}
\newtcblisting{risultato}{listing only, enhanced jigsaw, breakable,
  colback=black!5, colframe=black, boxrule=0.5pt, sharp corners,
  listing options={basicstyle=\ttfamily\small, numbers=none, language={},
    columns=flexible, keepspaces=true, breaklines=true}}

% --- Diagrammi di flusso (simboli standard) ---
\tikzset{
  terminale/.style = {ellipse, draw, minimum width=2.5cm, minimum height=0.9cm},
  io/.style        = {trapezium, trapezium left angle=70, trapezium right angle=110,
                      draw, minimum width=2.5cm, minimum height=0.9cm},
  processo/.style  = {rectangle, draw, minimum width=2.5cm, minimum height=0.9cm},
  decisione/.style = {diamond, draw, aspect=2, minimum width=2.5cm},
  freccia/.style   = {thick, -{Stealth}}
}

\title{Reti sequenziali e macchine a stati finiti}
\author{Marco Cerbella}
\date{Lezione 7}

\begin{document}
\maketitle

\section{Reti logiche sequenziali}

\subsection{Circuiti combinatori e circuiti sequenziali}

Nelle reti logiche combinatorie viste finora le uscite (il valore delle funzioni calcolate) dipendono \emph{esclusivamente} dalle configurazioni di bit in ingresso.

Nelle reti logiche \emph{sequenziali} invece la configurazione dei bit in ingresso non determina univocamente l'uscita: è necessario portarsi dietro un \emph{riassunto della storia del sistema} per determinare l'uscita.

\begin{esempio}[Riconoscitore del codice di Gray]
Un circuito sequenziale con due variabili binarie in ingresso $(x_0, x_1)$ e una variabile in uscita $z$, tale che $z = 1$ se e solo se viene riconosciuto il codice di Gray, cioè si riconosce una sequenza di configurazioni in ingresso nell'ordine
\[
(0,0) \to (0,1) \to (1,1) \to (1,0).
\]
\end{esempio}

\begin{esempio}[Il semaforo]
Si consideri il circuito controllore di un semaforo.
\begin{itemize}
    \item Il semaforo si può trovare in tre stati: rosso, giallo oppure verde.
    \item La lampada attualmente accesa dipende dalla temporizzazione associata a rosso, arancione e verde, e dallo stato precedente del semaforo.
\end{itemize}
\end{esempio}

\subsection{Parole chiave}

\begin{itemize}
    \item \textbf{Stato}: in pratica, lo stato codifica le informazioni relative alla storia del sistema utili a calcolare l'uscita del sistema.
    \item \textbf{Determinismo}: ci interessano sistemi deterministici: fissati l'ingresso e lo stato, l'uscita e l'evoluzione dello stato sono fissate.
    \item \textbf{Modello del tempo}: in una generica macchina a stati finiti (modello di un sistema fisico) l'evoluzione dello stato è dettata da un \emph{evento}. Se gli eventi sono prodotti da una sorgente indipendente di sincronizzazione (il segnale di \emph{clock}) si parla di macchine a stati \emph{sincrone}, altrimenti di macchine a stati \emph{asincrone}.
\end{itemize}

\section{Macchine a stati finiti (FSM)}

Le macchine a stati finiti sono un modello matematico per la descrizione di ``macchine'' dotate di memoria che eseguono algoritmi. Nel nostro contesto permettono di modellare ad alto livello i circuiti sequenziali (per esempio il controllore di un semaforo).
\begin{itemize}
    \item L'evoluzione del circuito è rappresentata come una sequenza di stati.
    \item Si è in uno \emph{stato presente} e si evolverà in uno \emph{stato futuro}.
    \item La realizzazione fisica consiste in una rete (combinatoria) per il calcolo delle uscite e dello stato futuro, a partire da ingressi e stato presente, e in un \emph{registro di attualizzazione} dello stato.
\end{itemize}

\begin{definizione}[Macchina a stati finiti: macchine di Mealy e di Moore]
Una macchina a stati finiti (sincrona) $M$ è una tupla
\[
M = \langle S,\, X,\, Z,\, \lambda,\, \delta,\, s_0\rangle
\]
dove:
\begin{itemize}
    \item $X = \{x_1, \dots, x_n\}$ è l'\emph{alfabeto di ingresso};
    \item $Z = \{z_1, \dots, z_m\}$ è l'\emph{alfabeto di uscita};
    \item $S = \{s_0, \dots, s_p\}$ è un insieme finito di \emph{stati} e $s_0$ è lo \emph{stato iniziale};
    \item $\delta$ e $\lambda$ sono due funzioni, rispettivamente di \emph{stato prossimo} e di \emph{uscita}:
    \begin{itemize}
        \item $z_k = \lambda(x_k, s_k)$ e $s_{k+1} = \delta(x_k, s_k)$ per le macchine di \textbf{Mealy};
        \item $z_k = \lambda(s_k)$ e $s_{k+1} = \delta(x_k, s_k)$ per le macchine di \textbf{Moore};
    \end{itemize}
\end{itemize}
dove $x_k \in X$, $z_k \in Z$ e $s_k \in S$ rappresentano rispettivamente il simbolo di ingresso, il simbolo di uscita e lo stato all'istante $k$.
\end{definizione}

Dunque in una macchina di Moore l'uscita dipende solo dallo stato; in una macchina di Mealy dipende dallo stato \emph{e} dall'ingresso.

\subsection{Rappresentazione con grafi di transizione dello stato}

I grafi di transizione dello stato sono una rappresentazione grafica molto comoda nelle fasi iniziali di progetto, in cui si passa da una descrizione formale della macchina sequenziale alla sua realizzazione fisica. Si tratta di un grafo orientato in cui ogni stato è rappresentato da un \emph{nodo} e ogni \emph{arco} corrisponde a una transizione dello stato. In particolare ogni transizione (arco) indica il prossimo stato in corrispondenza di un determinato ingresso.
\begin{itemize}
    \item In una FSM di \textbf{Moore} i simboli di uscita sono associati agli \emph{stati} (si scrivono dentro il nodo).
    \item In una FSM di \textbf{Mealy} i simboli di uscita sono associati alle \emph{transizioni} (si scrivono sull'arco, nella forma ingresso/uscita).
\end{itemize}

\tikzset{
  stato/.style = {circle, draw, minimum size=1.1cm, inner sep=1pt},
  arco/.style  = {-{Stealth}, thick}
}

\begin{esempio}[Tabella e grafo di una macchina di Moore]
Alfabeto di ingresso $\{0,1\}$, tre stati $s_0, s_1, s_2$ (iniziale $s_0$), uscita $U \in \{0,1\}$.
\begin{center}
\begin{tabular}{lccc}
\toprule
stato & \multicolumn{2}{c|}{ingresso} & uscita $U$\\
 & $0$ & $1$ & \\
\midrule
$s_0$ & $s_0$ & $s_1$ & $1$ \\
$s_1$ & $s_2$ & $s_1$ & $0$ \\
$s_2$ & $s_0$ & $s_0$ & $0$ \\
\bottomrule
\end{tabular}
\end{center}
\begin{center}
\begin{tikzpicture}[node distance=3cm]
  \node[stato] (s0) {$s_0/1$};
  \node[stato, right=of s0] (s1) {$s_1/0$};
  \node[stato, below right=1.6cm and 1.5cm of s0] (s2) {$s_2/0$};
  \path[arco]
    (s0) edge[loop above] node[above] {$0$} (s0)
    (s0) edge node[above] {$1$} (s1)
    (s1) edge[loop above] node[above] {$1$} (s1)
    (s1) edge[bend left=15] node[right] {$0$} (s2)
    (s2) edge[bend left=15] node[left] {$0,1$} (s0);
\end{tikzpicture}
\end{center}
\end{esempio}

\begin{esempio}[Tabella e grafo di una macchina di Mealy]
Stessi stati e alfabeto; ogni cella della tabella riporta \emph{stato prossimo/uscita}.
\begin{center}
\begin{tabular}{lcc}
\toprule
stato & \multicolumn{2}{c}{ingresso}\\
 & $0$ & $1$ \\
\midrule
$s_0$ & $s_0/1$ & $s_1/1$ \\
$s_1$ & $s_2/0$ & $s_1/0$ \\
$s_2$ & $s_0/0$ & $s_0/0$ \\
\bottomrule
\end{tabular}
\end{center}
\begin{center}
\begin{tikzpicture}[node distance=3cm]
  \node[stato] (s0) {$s_0$};
  \node[stato, right=of s0] (s1) {$s_1$};
  \node[stato, below right=1.6cm and 1.5cm of s0] (s2) {$s_2$};
  \path[arco]
    (s0) edge[loop above] node[above] {$0/1$} (s0)
    (s0) edge node[above] {$1/1$} (s1)
    (s1) edge[loop above] node[above] {$1/0$} (s1)
    (s1) edge[bend left=15] node[right] {$0/0$} (s2)
    (s2) edge[bend left=15] node[left] {$0/0,\ 1/0$} (s0);
\end{tikzpicture}
\end{center}
\end{esempio}

\section{Moore e Mealy a confronto}

Una macchina di Moore
\[
M_{\text{Moore}} = \langle S_{\text{Moore}},\, X,\, Z,\, \lambda_{\text{Moore}},\, \delta_{\text{Moore}},\, s_0\rangle
\]
e una macchina di Mealy
\[
M_{\text{Mealy}} = \langle S_{\text{Mealy}},\, X,\, Z,\, \lambda_{\text{Mealy}},\, \delta_{\text{Mealy}},\, s_0\rangle
\]
sono \emph{equivalenti}, dal punto di vista ingresso-uscita, se e solo se, ignorando l'uscita $z_0 = \lambda_{\text{Moore}}(s_0)$ di $M_{\text{Moore}}$, ogni possibile sequenza di simboli in ingresso genera una coppia di sequenze di simboli di uscita, una per $M_{\text{Moore}}$ e una per $M_{\text{Mealy}}$, identiche.

\subsection{Da Moore a Mealy}

È sufficiente associare l'uscita appartenente a uno stato a \emph{tutte le transizioni che partono da tale stato}.

\begin{center}
\begin{tikzpicture}
  \node[stato] (a) {$s/0$};
  \path[arco] (a) edge[loop above] node[above] {$1$} (a)
              (a) edge node[above, pos=0.9] {$0$} ++(-2,0);
  \node at (3,0) {$\leadsto$};
  \node[stato] (b) at (6,0) {$s$};
  \path[arco] (b) edge[loop above] node[above] {$1/0$} (b)
              (b) edge node[above, pos=0.9] {$0/0$} ++(-2,0);
\end{tikzpicture}
\end{center}

\subsection{Da Mealy a Moore}

\begin{itemize}
    \item Ogni stato della macchina di Mealy va \emph{replicato} tante volte quante sono le transizioni con uscite differenti verso tale stato.
    \item I nuovi stati avranno il valore di uscita uguale a quello della transizione da cui sono stati originati.
\end{itemize}
Per esempio, se a uno stato $s$ della macchina di Mealy arrivano transizioni con uscita $0$ e transizioni con uscita $1$, nella macchina di Moore $s$ diventa una coppia di stati $s/0$ e $s/1$: le transizioni con uscita $0$ arrivano in $s/0$, quelle con uscita $1$ arrivano in $s/1$, e da ciascuno dei due stati replicati partono le stesse transizioni che partivano da $s$.

\section{Esempi}

\subsection{Riconoscere la parola ``via''}

\begin{esercizio}
Si realizzi una macchina a stati finiti che riceve in ingresso una sequenza di caratteri e riconosce la parola ``via''. L'uscita prodotta è $0$ se non si sta riconoscendo la parola, $1$ mentre la si sta riconoscendo e $2$ una volta che sia stata riconosciuta.
\end{esercizio}

\begin{osservazione}
Nelle slide l'esercizio è solo enunciato. Una soluzione possibile (aggiunta qui, non presente nelle slide) è una macchina di Moore con quattro stati: $s_0$ (nessun carattere riconosciuto, uscita $0$), $s_1$ (visto ``v'', uscita $1$), $s_2$ (visto ``vi'', uscita $1$), $s_3$ (visto ``via'', uscita $2$). Da $s_0$: con ``v'' si va in $s_1$, con qualunque altro carattere si resta in $s_0$. Da $s_1$: con ``i'' si va in $s_2$, con ``v'' si resta in $s_1$, altrimenti si torna in $s_0$. Da $s_2$: con ``a'' si va in $s_3$, con ``v'' si torna in $s_1$, altrimenti in $s_0$. Da $s_3$: con ``v'' si va in $s_1$, altrimenti in $s_0$.
\end{osservazione}

\subsection{Il semaforo all'incrocio}

Si consideri un incrocio tra due strade, nord-sud (NS) ed est-ovest (EO), controllate da un semaforo.
\begin{itemize}
    \item Il semaforo può commutare ogni $30$ secondi: per semplicità si considerano solamente rosso e verde.
    \item È presente un sensore in grado di ``leggere'', per ogni direttrice, se esiste almeno un'auto in attesa oppure un'auto che si accinga ad attraversare (condizioni trattate allo stesso modo).
    \item Il semaforo deve cambiare colore, da rosso a verde, quando esiste un'auto in attesa sulla sua direttrice.
    \item Se ci sono auto in attesa su entrambe le direttrici il semaforo deve cambiare colore (al termine del tempo di commutazione).
\end{itemize}

\textbf{Ingressi.} Auto NS presente/assente: $i_{NS} = 1/0$. Auto EO presente/assente: $i_{EO} = 1/0$.

\textbf{Stati.} Semaforo NS verde (quindi semaforo EO rosso); semaforo EO verde (quindi semaforo NS rosso).

\textbf{Uscite.} Luce NS verde (quindi luce EO rossa); luce EO verde (quindi luce NS rossa). Nel nostro caso le uscite coincidono con gli stati.

Tabella di transizione (ingresso $(i_{EO}, i_{NS})$; uscita = luce verde dello stato):
\begin{center}
\begin{tabular}{lccccl}
\toprule
stato $S$ & $(0,0)$ & $(0,1)$ & $(1,0)$ & $(1,1)$ & uscita\\
\midrule
Verde NS & Verde NS & Verde NS & Verde EO & Verde EO & luce verde NS\\
Verde EO & Verde EO & Verde NS & Verde EO & Verde NS & luce verde EO\\
\bottomrule
\end{tabular}
\end{center}

\begin{center}
\begin{tikzpicture}[node distance=5cm]
  \node[stato, minimum size=2.2cm, align=center] (ns) {Verde NS\\\small luce verde NS};
  \node[stato, minimum size=2.2cm, align=center, right=of ns] (eo) {Verde EO\\\small luce verde EO};
  \path[arco]
    (ns) edge[loop left, looseness=6] node[left] {$00,\,01$} (ns)
    (eo) edge[loop right, looseness=6] node[right] {$00,\,10$} (eo)
    (ns) edge[bend left=20] node[above] {$10,\,11$} (eo)
    (eo) edge[bend left=20] node[below] {$01,\,11$} (ns);
\end{tikzpicture}
\end{center}
Le etichette sugli archi sono le coppie $(i_{EO}, i_{NS})$. La macchina è di Moore: l'uscita dipende solo dallo stato.

\section{Riferimenti ed esercizi per casa}

\begin{itemize}
    \item Capitolo~6 del libro di testo: C.~Bolchini et al., \emph{Reti Logiche}, Maggioli Editore. Tutto il capitolo tranne le sottosezioni 6.3.3, 6.3.4, 6.3.5.
    \item Esercizio~6.1 del libro di testo (le soluzioni sono sul libro di testo).
\end{itemize}

\end{document}
